\section{Concentrating the arity of a tensor network}\label{sec:local}


Our aim is to gather the large tables of a tensor network at a few
stations while preserving its contraction. The relevant budget is the
sum of the squared arities. For example, \(W/T\) stations of arity
\(T\) have budget \(WT\). Gathering \(D\) at a time into stations of
arity \(DT\) would multiply that budget by \(D\).
Lemma~\ref{lem:local} keeps the budget of the large stations at \(O(WT)\)
while reducing their number to \(O(W/(DT))\). Their total arity is then
\(O(W/\sqrt D)\). The transformation also introduces small tables.
Part~IV shortens the wires of this small-table network between successive
concentrations, then uses batched table lookup to evaluate the result.

A \emph{station} is a tensor given by an arbitrary table on an ordered
list of bits. Its arity counts the external bits as well as its ends of
internal wires. Group the stations into cells of a host graph of maximum
degree \(g\). Wires run within a cell, between adjacent cells, or to the
exterior. The same host may carry several networks. A cell used in two
networks counts twice. A station designated \emph{immobile} retains its
external bits and its cell. All external bits must belong to such stations.

\begin{lemma}[Local concentration]\label{lem:local}
Fix a finite field \(k\) with \(\omega_k<9/4\), a degree bound \(g\), and
an integer \(\ell\ge1\). For each sufficiently large fixed integer \(D\),
the following holds for all sufficiently large integers \(T\). Use an
extension \(K/k\) containing a fixed finite extension, with \(|K|\) at
least a sufficiently large polynomial in \(T\) and
\([K:k]=O_D(\log T)\). The input tables may already be over \(K\).

Suppose there are at most \(W/(DT)\) occupied cells and at most
\(W/(DT)\) immobiles in total. In any one network each cell has at most \(100D\) stations
of total arity at most \(100DT\), and over all stations of all networks
\(\sum_i t_i^2\le WT\). The tensor of each network is a sum over \(K\)
of networks with the same exposed bits. Every choice of one term per network, taken together, satisfies the following bounds.
\begin{enumerate}
\item There are \(O(W/(DT))\) \emph{frozen stations}, including the
holders of the original external bits in their original cells. Their
maximum arity is \(O(DT)\), their squared arities sum to \(O(WT)\),
and their arities sum to \(O(W/\sqrt D)\).
\item All other stations, called the \emph{core}, have arity at most
\(D^{-\ell}T\) and total arity \(O(W(1+\log D)^2)\).
\item A common host replaces each original cell by \(D^{O(1)}\) positions,
with possible contacts only within and between neighboring cells. The
induced host on the core positions has absolute bounded degree.
No two frozen stations touch by a wire. For some \(c=c(g,\ell,k)\) all active
internal bundles can have width at least \(\lceil T/D^c\rceil\), and
the bounds in the first two conclusions continue to hold when every
station is charged at least that arity.
\item The outer summation uses \(O_D(\log(T+1))\) control bits per
occupied cell. For its computational bound, represent each input table
by a label and its ports by ordered bundles, one per station pair,
specified by lengths and endpoint slots. Bits agree in order along a
bundle. Any local reordering is part of the supplied table label.
Given a fixed coloring with \(O_D(1)\) colors separating cells
within a bounded distance depending on \(D,g\), the validity of the indices, output
layout, and tracking of exposed sublists can be determined in \(O_D(1)\)
rounds of communication on the old host. With input labels of at most
\(A\) bits, the work and communication per occupied cell are
\(O_D(A+(1+\log T)^{O(1)})\) bit operations and bits. Each output table
has a label of length \(O_D(A+\log(T+1))\) in a dictionary fixed by the
input dictionary and the parameters.
\end{enumerate}
The implicit constants in the first three conclusions can depend on
\(k,g,\ell\), independently of \(D,T,W\). A fixed change of the numerical
input bound \(100\) changes only these constants.
\end{lemma}


The proof uses two sphere forms in each occupied cell.
The first joins neighboring cells along pairs of legs called \emph{seams}.
It represents a station by an insertion at a marked point. The energy
of this insertion bounds the square of the width needed to encode its
state. Moving insertions across seams and merging them at common marks
will concentrate their energy at the frozen stations. The seams have
a separate energy budget. A crossing spends some of this budget, and
at the end each seam is divided among many small tables.

To begin the motion we must realize the actual wires of the input
network. The second sphere, which stays in its cell, serves this purpose.
Corollary~\ref{cor:universal} allows its insertions to encode the
prescribed wiring on the first sphere. Both sphere legs participate in
mergers, but only the first can cross a seam. The two spheres use
different numbers of species and different energy scales. These choices
make the private leg left behind at a crossing thin.

\subsection{Sphere forms and a transposition identity}\label{subsec:seam}

We use the algebraic free-fermion construction of Section~\ref{sec:fermions}, now with many species.
Write \(V^{(r)}_e\) for the energy-\(e\) subspace of the exterior algebra
on generators \(\psi_{s,m}\) with species \(1\le s\le r\), mode \(m\ge0\),
and energy \(m+\tfrac12\). Use the wedge basis in lexicographic order
on \((s,m)\). Every basis vector has length congruent to \(2e\pmod2\). 

Put \(c_0=\pi/\sqrt3\). For \(q>0\) the single-species generating
function is \(Z(q)=\prod_{m\ge0}(1+\exp(-q(m+1/2)))\).
Since \(t\mapsto\ln(1+\exp(-t))\) is convex, integration on the
successive intervals of length \(q\) gives
\(\ln Z(q)\le q^{-1}\int_0^\infty\ln(1+\exp(-t))\,dt=\pi^2/(12q)\).
For \(r\) species, \(\dim V^{(r)}_e\le\exp(qe)Z(q)^r\).
Minimizing over \(q\) for \(e>0\), and using the vacuum for \(e=0\), gives
\begin{equation}\label{eq:local-width}
      \dim V^{(r)}_e\le \exp(c_0\sqrt{re}).
\end{equation}

We encode a leg by the binary index of its basis state, padding unused
codes with zero weights.

At distinct points \(z_i\in K\), the sphere form \(C_r\) sends a tuple
of wedges to a Pfaffian. List their generators in site order. Pairings
at a common site or between different species vanish. For a single
species the pairing of modes \(m,n\) at distinct sites \(i,j\), in that
order, is
\[
 [x^m y^n]\,(z_j+y-z_i-x)^{-1}.
\]
The Pfaffian of the empty matrix is one and an odd list gives zero.
Exchanging adjacent homogeneous sites of energies \(e,d\) costs the sign
\((-1)^{(2e)(2d)}\). In particular all such signs are scalar when the energies
are fixed. These definitions include characteristic two by reduction of
the integral alternating pairing formula.

If two adjacent sites in that order at \(z,z+u\) have energies \(e,m\), respectively, their
coefficient at \(u^{e'-e-m}\) is obtained by replacing the pair by one
site of energy \(e'\) at \(z\), precomposed with the map
\[
 I_{e'}:V^{(r)}_e\otimes V^{(r)}_m\longrightarrow V^{(r)}_{e'}.
\]
Only \(e'\ge0\) with \(e'-e-m\in\mathbb Z\) occur. To compute
\(I_{e'}\), sum over partial internal pairings between the two sites at
offsets \(0,1\), translate the unmatched generators by expanding around
zero, and wedge them in order at energy \(e'\). More precisely, the
translation at offset \(a\) replaces \(\psi_{s,d}\) by
\(\sum_{f\ge d}\binom fd a^{f-d}\psi_{s,f}\). This is the fusion
formula of Proposition~\ref{prop:fusion} extended to many species. Its proof decomposes a full
matching into its internal and external pairs and applies species by
species. As a function of any one position, the form has poles only at
the other positions and is bounded at infinity.

To join two spheres, contract the legs at two specified sites using
equality on the wedge basis at an energy \(e\). We call the pair a
\emph{seam}. To move an insertion through the seam, we need to transpose
a fusion coefficient against this ordinary equality pairing.

\begin{lemma}[Transposition at a seam]\label{lem:seam}
Let \(a,b\) be basis states at energies \(e,e'\), and
\(x\in V_m^{(r)}\). There are parity-preserving maps
\(S_j:V_m^{(r)}\to V_j^{(r)}\) for \(j\le m\), independent of \(a,b\),
and a nonzero scalar \(\alpha=\alpha(e,e',m)\) such that
\begin{equation}\label{eq:seam-adjunction}
 [I_{e'}(a,x)]_b
    =\alpha\sum_{j\le m}[I_e(b,S_jx)]_a .
\end{equation}
The coefficient is taken in the indicated wedge basis.
\end{lemma}
\begin{proof}
Include a square root \(i\) of \(-1\) in the fixed extension of \(k\).
Use \(y\) as the local coordinate at infinity, so the position is \(1/y\).
Before extracting a mode as a coefficient in \(y\), multiply the kernel
once for each generator there by
\(i/y\). The pairing between finite coordinate \(U\) and that chart is
then \(i/(1-Uy)\). Consequently the two-point form at zero and infinity
pairs equal basis states of length \(d\) with coefficient
\[
      (-1)^{d(d-1)/2}i^d=
      \begin{cases}1& d\ \text{even},\\ i& d\ \text{odd},\end{cases}
\]
and all other coefficients vanish. In the three-point form
at \(0,1,\infty\) on \(a,x,b\), first take the internal matchings at
zero and one. Pairing the remaining generators with infinity performs
exactly the translation and wedge defining \(I_{e'}(a,x)\). Thus the
three-point value is its \(b\)-coefficient times the displayed parity
factor.

Inverting the finite coordinate by \(V=1/U\) preserves the two-point
kernel when each argument receives the factor \(-i/U\). Indeed
\((-i/U_1)(-i/U_2)/(1/U_2-1/U_1)=1/(U_2-U_1)\).
The old infinity chart becomes an ordinary finite chart with factor one,
and the old zero becomes an infinity chart with one minus sign per
generator. At the middle site, take the exterior power of the
single-generator substitution
\[
 \psi_{s,d}\longmapsto
 \sum_{f=0}^{d}
 \left([X^d]\,\frac{-i}{1+X}
       \left(\frac{1}{1+X}-1\right)^{\!f}\right)\psi_{s,f}.
\]
Projection to energy \(j\) gives \(S_j\). Each term retains the length
and has energy at most \(m\). Reverse the site order and evaluate the
three-point form with \(b\) at zero and \(a\) at infinity. The order
reversal, the signs from the old zero, and the two diagonal parity factors
depend only on the three original parities. Taking their ratio proves
\eqref{eq:seam-adjunction}.
\end{proof}

For one species take \(a=1\) and \(b=x=\psi_0\). At energy
\(1/2\) the near-side fusion passes \(x\) unchanged. After inversion the
middle insertion is \(-i\psi_0\), and its far-side contraction with
\(b\) at offsets \(0,1\) gives \(-i\) at energy zero. The factor
\(\alpha=i\) recovers the original coefficient one. The same exchange
of input and output levels will carry a general insertion across a seam.

\section{Planting the wiring at marked points}\label{subsec:planting}

A station in the input may have wires to many other stations. To move it
as one insertion, we encode its ends of those wires in a single leg.
The sphere form must then test whether the bits at opposite ends agree.
We implement that test separately in each occupied cell of each network.

There are two sphere forms in a cell. The form $A$ carries the station
legs and the seams to neighboring cells. Extra marks on $A$ supply a
joint state that tests the wiring. The private form $B$ realizes this
joint state through a separate insertion table at each extra mark.
Thus, after planting, each insertion has just one leg on each sphere,
along with any exposed or control bits. At original stations and seams
the $B$ leg is the vacuum. Figure~\ref{fig:stack} shows the two kinds
of insertions.

\begin{figure}[htb]\centering
\begin{tikzpicture}[x=1cm,y=1cm,font=\small,
  strip/.style={draw=ink,rounded corners=3pt,line width=.65pt},
  ins/.style={draw=ink,rounded corners=2pt,fill=white,inner xsep=7pt,inner ysep=4pt},
  leg/.style={draw=ink,line width=.8pt}]
\draw[strip,fill=ink!10] (-4.35,.98) rectangle (4.35,1.55);
\draw[strip,fill=ink!5] (-4.35,-1.55) rectangle (4.35,-.98);
\node[anchor=west] at (-4.22,1.26) {Sewn sphere \(A\)};
\node[anchor=east] at (4.22,1.26) {\(R\) species};
\node[anchor=west] at (-4.22,-1.26) {Private sphere \(B\)};
\node[anchor=east] at (4.22,-1.26) {\(R/b\) species};
\draw[leg] (-5.35,1.26)--(-4.35,1.26);
\draw[leg] (4.35,1.26)--(5.35,1.26);
\node[above,font=\footnotesize] at (-4.85,1.26) {seam};
\node[above,font=\footnotesize] at (4.85,1.26) {seam};
\node[ins] (sta) at (-2.15,0) {original station};
\node[ins] (aux) at (2.0,0) {auxiliary insertion};
\draw[leg] (sta.north)--(-2.15,.98);
\draw[leg,densely dotted] (sta.south)--(-2.15,-.98);
\node[anchor=west,font=\footnotesize] at (-2.1,-.6) {vacuum};
\draw[leg] (aux.north)--(2.0,.98);
\draw[leg] (aux.south)--(2.0,-.98);
\node[anchor=west] at (2.05,.62) {\(z^P\)};
\node[anchor=west] at (2.05,-.62) {\(z\)};
\foreach \x/\y in {-2.15/.98,2.0/.98,2.0/-.98}
    \fill[ink] (\x,\y) circle (1.5pt);
\draw[draw=ink,fill=white] (-2.15,-.98) circle (1.5pt);
\end{tikzpicture}
\caption{The station table encodes its incident internal wire bits on \(A\).
The auxiliary tables couple \(A\) to \(B\) and together impose equality
along the wires. Only \(A\) is sewn to neighboring cells.}\label{fig:stack}
\end{figure}

A mark at $z$ on $B$ corresponds to the mark at $z^P$ on $A$, where $P$
is a power of the characteristic. Moving $z$ by $t$ moves $z^P$ by $t^P$.
In units of the Laurent order on $A$, the energy of the pair is therefore
$e_A+P^{-1}e_B$. A given energy budget allows $P$ times as much energy
on $B$. We keep the corresponding leg narrow by using fewer species on
$B$. This will bound the width left behind when a station crosses a seam.

Fix $g,\ell$ as in Lemma~\ref{lem:local}. Let $F$ be the first
characteristic power at least $D^{1000(\ell+1)}$ and let $a=\lceil F^3\rceil$.
For a chunk count $M$ and a sufficiently large characteristic power $b>F$,
take $R$ to be the largest multiple of $bM$ at most $T/a$. Set
\begin{equation}\label{eq:plant-scales}
 R'=R/b,\qquad P=b/F,\qquad \kappa=P^{-1},\qquad H=Ra^2 .
\end{equation}
Here $R$ and $R'$ count the species on $A$ and $B$, respectively, and
$T/(2a)\le R\le T/a$ for sufficiently large $T$. If $e_A+\kappa e_B\le h$,
Equation~\eqref{eq:local-width} bounds the widths by constant multiples
of $\sqrt{Rh}$ and $\sqrt{R'Ph}=\sqrt{Rh/F}$. In bits, up to rounding,
\begin{equation}\label{eq:plant-widths}
 w_A=O\bigl(T\sqrt{h/H}\bigr),\qquad
 w_B=O\bigl(T\sqrt{h/(FH)}\bigr).
\end{equation}

\subsection{The wiring form for one species}

Combine the wires between each pair of distinct stations into one ordered
bundle, absorbing endpoint permutations into the station tables. Contract
self-loops in those tables first. Split each internal bundle into packets
of $R$ bits, padding the last packet with zeros. Species $s$ will carry
bit $s$ of every packet. A station meets $O_g(D)$ bundles, and a cell
together with its incident edges meets $O_g(D^2)$. At a station of
arity $t_i$, one species therefore carries at most $t_i/R+O_g(D)$ bits.

Choose distinct wedge basis states at an integer energy $h_i$ to encode
the bit strings at station $i$. The logarithm of the single-species
dimension is asymptotic to $(\pi/\sqrt3)\sqrt{h_i}$ through integer
levels by Lemma~\ref{lem:dimension}. Since $R\ge T/(2a)$ and $a\ge D$,
we may take
\[
                 h_i=O_g\bigl(a^2(1+t_i/T)^2\bigr).
\]
At seam $j$ the packets across the interface contribute
$O_g(DT/R+D^2)=O_g(Da)$ bits to one species. An integer energy
$q_j=O_g(D^2a^2)$ suffices. The two cells use the same bit order and
the same basis codes there. For fixed $D$ the packet counts per
species, energies and codes range over finite sets.

Define the multilinear form $L$ on these station and seam levels by
its coefficients. Give a basis tuple coefficient one if all codes are
valid and both ends of every bit agree. For a wire to another cell,
compare its station end with its seam end. Every other coefficient
is zero. We will sew a multispecies seam at a fixed \emph{total}
energy. To recover the correct energy in each individual species,
put $Q=\sum_j q_j$ and extend $L$ by zero to all other tuples of
nonnegative half-integer seam energies of sum at most $Q$. The station
energies stay fixed.

\begin{lemma}[Auxiliary marks of bounded weight]\label{lem:plant-probes}
Inserting a joint state at polynomially many extra marks on a
single-species sphere realizes the extended form $L$. At each extra
mark the state uses only the vacuum and $\psi_0$. In every word in its
support the number of occupied marks is even and at most $O_g(Da)$.
\end{lemma}
\begin{proof}
A wedge of length $d$ has energy at least
$\tfrac12+\tfrac32+\cdots+(d-\tfrac12)=d^2/2$. At the stations the
total length is $O_g(a\sum_i(1+t_i/T))=O_g(Da)$. At the at most $g$
seams the same bound follows from Cauchy--Schwarz and their total
energy bound $Q=O_g(D^2a^2)$. List all pairs $(v,m)$ of original
marks and modes occurring in the resulting input spaces. The list
has polynomial size in $D,a$.

An insertion of $\psi_0$ at a new point $z$ pairs with those modes by
the covector with coordinates $(z-z_v)^{-m-1}$. Distinct principal
parts give linearly independent rational functions in every
characteristic. Choose distinct regular points whose covectors span
the dual mode space. Indeed, a nonzero linear combination annihilated
by the points already chosen is nonzero at some unused point of the
algebraic closure. At most one point per listed mode is needed, all
in a finite extension. Add vacua if necessary to obtain at least
three auxiliary marks.

For a subset of $d$ new marks, insert $\psi_0$ there and the vacuum
at the others. Consider first the part of the resulting Pfaffian
in which every auxiliary generator pairs with an original generator.
This part is itself a linear combination of insertion tests on at
most $d$ marks. Indeed, order all auxiliary generators last and group
the matchings by their exact set of internal auxiliary pairs. Each
nonempty set contributes its contracted scalar, with a sign depending
only on those paired positions, times the part on the remaining marks
with no internal auxiliary pairs. Subtracting these contributions
inductively isolates the desired part.

Now take a nonzero vector in the even part of the direct sum of input
spaces, and let $d$ be its greatest nonzero total wedge length. A test
on $d$ marks with no internal auxiliary pairs vanishes on every shorter
wedge. On the length-$d$ component it evaluates the exterior product of
the chosen $d$ covectors, up to a common sign. The original basis
tuples are distinct basis wedges in the exterior algebra on the list
of pairs $(v,m)$. Some choice of $d$ covectors therefore detects the
component. The resulting even tests span the full dual of the even
part. Each actual insertion used has even weight at most $d=O_g(Da)$
and vanishes on odd vectors. The prescribed integer energies give
even lengths, so the zero extension of $L$ also vanishes on the odd
part. A linear combination of insertion tests now gives the required
joint state.
\end{proof}

Use one copy of this state in each of the $R$ species. At station $i$
the insertion table enforces level $h_i$ in every species, decodes the
packets, checks that padded bits are zero and supplies the original
table entry, with its exposed bits. At seam $j$ we sew by basis
equality at total energy $Rq_j$. To see why this recovers the original
wiring, group the Pfaffian by species. Order the original and auxiliary
marks in two consecutive lists and reorder each list by species.
The auxiliary reordering sign can be absorbed in the joint state.
Interleaving the lists costs no further sign because every species
has even auxiliary weight. The original reordering sign depends only on
the input wedges. In any nonzero product after the auxiliary state has
been contracted, the seam tuple of each species
is either the prescribed tuple or has sum greater than $Q$, by the zero
extension. The sums over all species add to $RQ$, so every species
has the prescribed tuple. All the original wedges then have even
length and the remaining reordering sign is one. The product tests
exactly the equality of the bits along each wire. We also place one
vacuum collector in each cell for the motion step.

\subsection{Realizing the joint auxiliary state}\label{subsec:plant-private}

Lemma~\ref{lem:plant-probes} prescribes one tensor shared by all
auxiliary marks. For motion we need a separate insertion table at
each mark. The private sphere supplies the shared tensor through
local maps. One $B$ species will serve a block of $b$ species on $A$.

Partition the $R$ sewn species into $R'$ consecutive blocks. A block
has \emph{type} $\tau=(d_i)_{i=1}^m$ if $d_i$ of its modes $\psi_0$
are occupied at auxiliary mark $i$. Projecting the block tensor onto
these local subspaces gives a tensor $J_\tau$, including the sign
for reordering within the block. Its $i$th leg has dimension
$\binom b{d_i}$. For a fixed ordered word of block types, the joint
tensor is the product of the $J_\tau$ up to a known sign. Exchanging
two site lists from different blocks contributes only the product of
their lengths modulo two to the sign exponent. Retain the types with
$J_\tau\ne0$.

Write $H_2$ for binary entropy in bits, with $H_2(0)=H_2(1)=0$.
Round each $H_2(d_i/b)$ up to a strictly larger point $u_i$ of a
fixed positive grid of mesh $1/m$. Then
\begin{equation}\label{eq:plant-ent}
 \log_2\binom b{d_i}\le b u_i,\qquad
 \max_i u_i\le 2,\qquad \sum_i u_i=O_{g,\ell}(Da\log D).
\end{equation}
Indeed, Lemma~\ref{lem:plant-probes} gives $s=\sum_i d_i/b=O_g(Da)$.
Concavity bounds the sum of the entropies by $s\log_2(em/s)$,
with value zero at $s=0$. The $m$ marks are polynomially many in $D$
for fixed $\ell$ and characteristic, and rounding adds at most one
to the sum.

At the corresponding $B$ marks, take the direct sum $S_\tau$ of all
nonzero single-species sphere forms with energy at most $C(bu_i)^2$
at site $i$, where $C$ is an absolute constant chosen below. For
sufficiently large $b$ there is a constant $c_\tau$ such that for
every $d\ge1$ a product of local linear maps gives
\begin{equation}\label{eq:plant-comparison}
 \bigoplus_{s=1}^{c_\tau} S_\tau^{\otimes d}
                    \longrightarrow J_\tau^{\otimes d}.
\end{equation}
The independence of $c_\tau$ from $d$ will let one selector index
serve every occurrence of the type.

For the comparison, take natural logarithms and work over the
algebraic closure, where the exponent is still strictly below $9/4$.
For a normalized spectral price at $x_1,\ldots,x_m$ put
$M_u=\inf_o\sum_i u_i\,d(o,x_i)$. Preparing $J_\tau$ at $o$ and
transporting its legs bounds its price by $b(\log 2)\sum_i u_i d(o,x_i)$
by Equation~\eqref{eq:plant-ent}. Taking the infimum gives $b(\log 2)M_u$. Applying the same transport bound to
the size-two diagonal, whose price is $\log 2$, gives
$\inf_o\sum_i d(o,x_i)\ge1$. Hence $M_u\ge\min_i u_i>0$.
Apply Corollary~\ref{cor:universal} at scale $N=b$ with width caps
$c_0\sqrt C\,u_i$. Its comparison tensor is precisely $S_\tau$,
whose price is $b(c_0\sqrt C\,M_u+o(1))$ uniformly in the normalized
price.
Choose $c_0\sqrt C>2\log 2$. For fixed $D$ there are finitely many
grid vectors and choices of $A$ positions, all defined over a fixed
finite extension. Taking $P$th roots permutes that field, so the
possible $B$ positions also range over a finite set as $b$ varies.
It follows that $S_\tau$ strictly dominates $J_\tau$ in every price
for all sufficiently large characteristic powers $b$.
The field-extension and strict comparison results of
Section~\ref{sec:prices-background} give Equation~\eqref{eq:plant-comparison}
over the fixed field of definition. The maps extend to the field of
the input tables.

Fix the type word. For each type of multiplicity $d>0$, expand the
restriction in Equation~\eqref{eq:plant-comparison} by choosing one
copy index $s$ and one $B$-energy tuple for each of the $d$ factors.
Assign the factors to private species in the original block order.
Their product is the multispecies sphere form up to the sign for
permuting the homogeneous lists. Each site map now supplies a local
table coupling the $A$ and $B$ legs at that mark. It may act jointly
on all blocks of a given type \emph{at that mark}. Summing the choices
and type words produces the required auxiliary tensor. The $B$ legs
at the original stations and seams stay in the vacuum.

\paragraph{Control and storage.}
A \emph{private selector} records the ordered type word, the
per-block $B$ energies and one index $s$ per occurring type. A single
table checks validity, including zero padding of unused slots, and
supplies the permutation sign. The types and twice the energies cost
$O(m\log(2b))$ bits per block, since $d_i\le b$ and the doubled
energies are $O(b^2)$. The copy indices cost $O_{D,b}(1)$ bits in
total. With $R/b$ blocks, choosing first $b$ and then $T$ large makes
the selector shorter than $\epsilon T$ for any fixed $\epsilon>0$.

On $A$ every occupied auxiliary mode costs energy $1/2$, giving
total auxiliary energy $O_g(RDa)$. On one $B$ species the sum of
the caps is $O(b^2\sum_i u_i^2)=O_{g,\ell}(b^2 Da\log D)$ by
Equation~\eqref{eq:plant-ent}. Multiplying by $R'$ and converting to
$A$ units gives the cell bound
\[
 \sum_i(e_{A,i}+\kappa e_{B,i})
    =O_{g,\ell}(RDa+F RDa\log D)\le D H
\]
for sufficiently large $D$, since $H=Ra^2$ and $a\ge F^3$.
For storage, use the actual spaces of the selected block. At mark $i$
either leg needs only $O(bu_i+1)$ bits, by
Equations~\eqref{eq:plant-ent} and~\eqref{eq:local-width}.
Over all blocks, each auxiliary insertion needs $O(R)=O(T/a)$ bits
for these two legs. Summed over the marks, this costs
$O_{g,\ell}(DT\log D)$, taking $b\ge m$ to absorb rounding.

The selector is summed inside the network, but the outer layout must
know the energy and width at each mark. Split the blocks into $M$ equal
consecutive chunks. In the outer sum, specify for each chunk the
histogram of its pairs (type, $B$-energy tuple). A histogram records
the multiplicities of the pairs, while the private selector records
their order. The alphabet is fixed independently of $T$, so the
histograms cost $O_{D,b,M}(\log(T+1))$ bits per cell. They fix the
energy and the concatenated code length of each leg in each chunk.
The table checks the histograms against the selector and uses its block
order to decode. Each valid selector occurs in exactly one outer term.
Thus $M$ can be chosen for the chunk layout, then $\epsilon$ for its
selector routing, and finally $b$ for the comparison and selector width,
all before $T$ grows.

\section{Moving insertions through a sewn surface}\label{subsec:motion}

The planting construction expresses the original tables on pairs of
spheres in each cell. We now combine these tables at a small number
of marks. Two outcomes of a Laurent expansion suffice. At an insertion
mark, a coefficient fuses the moving table with the table already
there. At a seam, a sufficiently negative coefficient transfers the
moving table to the neighboring cell and lowers the energy of the
seam. This decrease bounds both the number of crossings and the amount
of energy that any cell can receive.

In each cell, designate the marks of the original immobile stations
and one additional vacuum mark as \emph{collectors}. They stay in their
cell. Process the original movable stations in turn, skipping those
already absorbed. The moving table keeps its ordinary bit ports. If its
current legs on the sewn sphere $A$ and private sphere $B$ have
energies $m_A,m_B$, put
\[
       m=m_A+m_B/P,\qquad \mu=m+H,\qquad L=\lceil 2P\mu\rceil .
\]
Here $P$ is the exchange rate of Equation~\eqref{eq:plant-scales}.
The reserve $H$ stays with this move from its start until it ends, while
$m$ is updated after every fusion or crossing.

\subsection{Reconstructing an evaluation from collisions}

Fix the branch of the outer sum and, for the following identity, the
private selectors. The branch determines the energies, so the same
identity works for all allowed selector values. Hold all other positions
and insertion tables fixed. Each entry of the contraction is a rational
function of the mover's $B$-coordinate $z$. Its $A$-coordinate is $z^P$.
The sphere kernels have poles only at other marks on the current sphere
and are bounded at infinity. Near a mark with $B$-coordinate $z_0$, write
$s=z-z_0$. Since $P$ is a power of the characteristic, the corresponding
coordinate difference on $A$ is $s^P$. Hence an $A$-fusion of order $i$
and a $B$-fusion of order $v$ contribute to order $j=Pi+v$ of the
contraction. At an insertion of weighted energy $m_0$ the merged energy is
\begin{equation}
            m_{\rm new}=m_0+m+j/P.                     \label{eq:motion-fusion}
\end{equation}
Indeed each fusion adds its order to the two incident energies, by
Proposition~\ref{prop:fusion}. Nonnegativity of the two output energies
bounds $i$ and $v$ below. Thus the pole orders are bounded, and only
finitely many pairs $(i,v)$ contribute at each $j$.

To reconstruct the value at the current position, use the following
coefficients at the other marks of the cell.
\begin{itemize}
\item at every insertion, its entire principal part;
\item at the designated vacuum collector, also the orders $0,\ldots,gL$,
even after this collector has absorbed a table;
\item at each seam, only the orders $j\le -L$ within its pole bound.
\end{itemize}
A function on which all these coefficients vanish has poles only at the
at most $g$ seams, of order at most $L-1$ each. It has a zero of order
at least $gL+1$ at the designated collector and no pole at infinity.
The degree of the divisor of a nonzero rational function on the projective
line is zero, so this function vanishes identically. On the space with
the prescribed pole bounds, evaluation at the mover's position is
therefore a linear combination of the listed coefficients. This
combination acts entrywise, independently of the tables.

A coefficient at an insertion is a sum over the compatible pairs
$(i,v)$. For each pair, apply the two fusion maps and contract them
with the moving and stationary tables, including the sign of their
chosen order. A collision with a collector ends the move. At any
other insertion only $j<0$ is used, so the combined table continues
with at most the sum of the two original weighted energies. This step
can absorb an auxiliary insertion, an original station awaiting its
turn, or a private insertion left by an earlier crossing.

\subsection{Crossing a seam}

Suppose the selected coefficient has order $j\le-L$ at a seam of
energy $e$. Only sphere $A$ is sewn there. The insertion on $B$ at the
seam is the vacuum, so the $B$-expansion has Taylor orders $v\ge0$.
For each such $v$ with $i=(j-v)/P$ integral and in the fusion range,
the new $A$-energy at the seam satisfies
\begin{equation}
       e'=e+m_A+i\le e+m_A-2\mu\le e-\mu .             \label{eq:motion-drop}
\end{equation}
Terms with $e'<0$ vanish. With $i$ fixed, interpolate the $B$ Taylor
coefficient using evaluations at fresh regular positions. Each resulting
term has one private insertion of energy $m_B$ left behind, with
vacuum on $A$. Its index is carried to the original table by an ordinary
equality wire. We call this private insertion the \emph{residue} of the
crossing.

To identify the new moving table, fix its old $A$-state $x$ by
linearity. Let $a,b$ range over wedge basis states of energies $e,e'$,
respectively, and denote the two remaining sphere forms by
$C_{\rm near},C_{\rm far}$, suppressing their spectator arguments.
Lemma~\ref{lem:seam} gives the identity
\[
\begin{aligned}
 &\sum_b C_{\rm near}(b)\sum_a [I_{e'}(a,x)]_b\,C_{\rm far}(a)\\
 &\hspace{1em}=\alpha\sum_{t\le m_A}\sum_b C_{\rm near}(b)
                      \sum_a [I_e(b,S_t x)]_a\,C_{\rm far}(a).
\end{aligned}
\]
On the right, the inner sum is the Laurent coefficient of order
$e-e'-t$ for an insertion $S_t x$ approaching $b$ on the far-side
$A$-sphere. This order is integral, since $e'-e-m_A=i$ and $m_A-t$
are integers for nonzero terms. Interpolate the coefficient by far-side
regular evaluations, with vacuum in the corresponding $B$-position.
The sum over $b$ sews the two spheres at energy $e'$. The new moving
table, now on the far side with energies $(t,0)$, retains its ordinary
ports and the wire to the residue (Figure~\ref{fig:motion-crossing}).
All reordering signs depend only on the specified energies, and the
identity remains valid when $x$ is coupled to the private index or to a
spectator.

\begin{figure}[!htbp]
\centering
\begin{tikzpicture}[x=1cm,y=1cm, font=\small,
 cell/.style={draw=ink!50,fill=ink!3,rounded corners=2pt},
 tab/.style={draw=ink,fill=white,rounded corners=1pt,inner sep=3.2pt,
             font=\small,align=center}]
 \node[font=\small\bfseries] at (-3.62,1.45) {Near cell};
 \node[font=\small\bfseries] at (3.62,1.45) {Far cell};
 \draw[cell] (-5.36,-1.11) rectangle (-1.88,1.13);
 \draw[cell] ( 1.88,-1.11) rectangle ( 5.36,1.13);
 \node[align=center] at (-3.62,.55) {sewn sphere $A$};
 \node[align=center] at ( 3.62,.55) {sewn sphere $A$};
 \fill[sea] (-1.88,.55) circle (2pt);
 \fill[sea] ( 1.88,.55) circle (2pt);
 \draw[sea,line width=.9pt] (-1.88,.55) -- (1.88,.55);
 \node[fill=white,inner sep=2pt,text=sea] at (0,.55) {seam energy $e'$};
 \node[tab] (res) at (-3.62,-.40) {private residue\\ $(0,m_B)$};
 \node[tab] (mov) at ( 3.62,-.40) {moving table\\ $(t,0)$};
 \draw[dashed,rust,line width=.75pt] (res.east)--(mov.west);
 \node[fill=white,inner sep=2pt,text=rust] at (0,-.40) {$B$-state index};
\end{tikzpicture}
\caption{A moving insertion of energies $(m_A,m_B)$ crosses from the near
cell to the far cell. The solid wire identifies seam states of energy $e'\le e-\mu$. The dashed
wire carries the old private state to the new moving table, whose sewn
energy is $t\le m_A$.}
\label{fig:motion-crossing}
\end{figure}

Both interpolations can be made over the same finite field $K$ as the
original contraction. Multiplying by a common denominator embeds the
space of rational functions bounded at infinity with prescribed poles
into the polynomials of degree at most the total pole bound. Evaluation
at more than that many regular field points is injective, so these
evaluations span the dual and recover each required Laurent coefficient.
Exclude all occupied positions. The map $z\mapsto z^P$ permutes $K$, so
regular evaluations can use corresponding positions on both spheres.
The accounting below bounds the degrees by $O_D(T)$ and the number of
occupied positions per cell by a polynomial in $D$. One can first
obtain those bounds over the algebraic closure and use them to choose
$K$ before any interpolation. Fix an ordering of the field
elements and, for each allowed set of regular points in this order,
store the entries of one reconstruction matrix in the table dictionary. The local control need only
name the selected point and check its availability.

Allow at most $U=D^5$ crossings in a move. When this limit is reached,
its current mark becomes a collector. Nonterminal fusions between
successive crossings remove one stationary insertion each. A completed
move never releases a new mover, so this procedure terminates in every
summand.

\subsection{Accounting for motion}

We first bound the energy stored in the final insertions. An original
station of arity $t_i$ has initial weighted energy $O((1+t_i/T)^2H)$.
The auxiliaries contribute at most $DH$ per occupied cell. Also allocate
one reserve $H$ to each original movable station. The initial insertion
energy and reserves are $O(D^2H)$ per cell. Over all cells, including
all networks in the batch, the bound is
\[
 O\!\left(H\left(D\,\#\{\text{occupied cells}\}
                  +\#\{i\}+T^{-2}\sum_i t_i^2\right)\right)
       =O(WH/T).
\]

A continuing fusion cannot increase the sum of participating energies,
by Equation~\eqref{eq:motion-fusion}. Keep only the reserve of the
current move. In a crossing the energy $m$ pays for $t\le m_A$ on the
far side and $m_B/P$ at the near-side residue. At a final fusion the
collector may gain up to $m+gL/P\le(1+3g)\mu$, since $H\ge1$ and
$P\ge1$. Count energy at collectors with weight $1/(1+3g)$ and energy
elsewhere, including reserves, with weight one. The sum cannot
increase at any of the three operations or when a move expires. The
total insertion energy in every resulting summand is consequently
$O_g(WH/T)$.

There is also a bound in each cell. A crossing brings in at most
$t+H\le\mu$ of charge. By Equation~\eqref{eq:motion-drop} the seam loses
at least this much energy. A cell has at most $g$ seams, each initially
of energy $O_g(D^2H)$. Adding the initial charge within the cell bounds
its current insertion energy by $O_g(D^2H)$ at every step. Moreover,
each crossing decreases the seam energy by at least $H$, so throughout
the construction a cell has $O_g(D^2)$ crossing incidences. Only these
incidences can introduce new stationary residues. With $H=O_D(T)$ and
$P$ fixed, these energy and mark bounds give the asserted $O_D(T)$ total
pole bound. The number of expired moves over all cells is at most the
total number of crossings divided by $U$. Including the initial immobiles
and designated collectors, the total number of collectors is $O(W/(DT))$.

It remains to route the ordinary wires created in the expansion. By
Equation~\eqref{eq:plant-widths}, the private state left at a crossing
uses $O_g(DT/\sqrt F+1)$ bits. One endpoint follows the mover. The
other stays at the residue until another insertion absorbs it, then
travels with that insertion. A residue that can still be absorbed
carries only the endpoint created at its own crossing. Once a moving
endpoint finishes a move, its collector stays fixed. A move visits at
most $U+1$ cells counting repetitions, collecting at most $O_g(D^2)$
residue endpoints per visit and creating at most $U$ of its own.
Hence it carries $O_g((U+1)D^2)$ endpoints along any one hop. Using
the bound on incident crossings and including terminations gives the
looser bound $O_g((U+1)D^6)$ local wire segments per cell. Successive
motions within a single cell require only one local connection. With
$F$ from Section~\ref{subsec:planting}, these segments have total width
$O(D^{-2\ell-1}T)+O_D(1)$ there.

The private-selector wires start at marks that read them and at the
stations checking the selectors and evaluating the sphere forms. Their
endpoints follow the same paths. Each cell starts with polynomially
many such ports in $D$ and the number of chunks, so the path and
crossing bounds give a polynomial bound in the same quantities on
incident selector segments and copies. This count is independent of
$b$ and of the selector rate. After fixing the number of chunks, choose
the rate small enough that the total selector width, including equality
copies, is at most $D^{-2\ell}T$ per cell. The next section converts the
remaining energy levels into table arities.


\section{Accounting for the new arities}\label{subsec:local-output}


After the motion, each collector has two homogeneous sphere legs and may
also hold external bits or thin wires. A seam can still have large energy.
We divide the species into groups so that each group meets a small table,
even when one group receives the entire energy of a seam.

Choose \(M\), a power of \(D\) specified below. Both species counts
\(R,R'\) can be taken divisible by \(M\) in
Equation~\eqref{eq:plant-scales}. Divide the species of each sphere
into \(M\) consecutive equal chunks, preserving whole blocks in the
private construction. Expand each remaining homogeneous leg as a
direct sum over its tuple of chunk energies, using the prescribed tuple
at each untouched auxiliary insertion and the same tuple on both sides
of a seam. Every basis state occurs once in this expansion. For fixed
chunk energies, the sphere form factors into the chunk forms.
Reordering the generators contributes a single sign determined by the
parity of each chunk at each site.

Consider a generic leg of energy \(e\) on \(r\) species, and write
\(e_j\) for the energy in chunk \(j\), so \(\sum_j e_j=e\).
By Equation~\eqref{eq:local-width} and Cauchy--Schwarz, the total
coded width satisfies
\begin{equation}\label{eq:local-chunk-size}
 \sum_{j=1}^M\left\lceil\frac{c_0}{\ln 2}\sqrt{(r/M)e_j}\right\rceil
 \le \frac{c_0}{\ln 2}\sqrt{r e}+M .
\end{equation}
A layout can use the integer width \(\lceil4\sqrt{(r/M)e_j}\rceil\)
instead. This amounts to a rounded integer square root since \(2e_j\)
is an integer, and the same bound holds with coefficient four.

An untouched auxiliary
leg uses the actual product subspace specified by its selector. Within a
chunk, the total number of bits can be read from the counts of block types
and block energies, since a block of type \((d_i)_i\) with rounded entropy bounds
\(u_i\) needs at most \(O(bu_i)\) bits on leg \(i\). Under a fixed selector these codes inject
into the corresponding tensor product of homogeneous spaces. The table
for the chunk form decodes them and assigns weight zero to unused words.

One arbitrary table implements the form at each cell and chunk on the
first and, separately, the private sphere. The number of legs in a cell
is polynomial in \(D\). Equation~\eqref{eq:local-chunk-size}, applied
to a single chunk, bounds a generic port there by \(O(DT/\sqrt M+1)\),
and the bounds on actual auxiliary subspaces give \(O(R/M+1)\) per
auxiliary leg and chunk. Thus the total arity at one of these tables is
at most \(D^{c_1}T/\sqrt M+O_D(1)\), before the small selector wires
are supplied. The exponent \(c_1\) can be chosen independently of
\(M,b\) and the selector rate. Take \(M\) to be a sufficiently large
power of \(D\) to put this quantity below \(D^{-\ell}T/10\). Use
selector rate small enough for the routing bounds of
Section~\ref{subsec:motion} with this \(M\), then choose the block
size \(b\) in the planting. This order permits all three choices to be
made before \(T\) grows.

\paragraph{Frozen stations.}
The collectors become the frozen stations. At one collector of weighted
energy \(e\), the total width of its two sphere legs is
\(O(T\sqrt{e/H})+O(M)\) by Equations~\eqref{eq:plant-widths}
and~\eqref{eq:local-chunk-size}. The motion bounds give
\(\sum e=O(WH/T)\) over all collectors and \(e=O(D^2H)\) at any
one of them. Thus the energy term in the width has squared sum \(O(WT)\)
and maximum \(O(DT)\). The sum over collectors of the squares of their total external widths
is controlled by the input bound, since every immobile retains its own
external ports. All additional wires from motion and selectors together
have total width \(O(D^{-2\ell}T)\) per cell. Their squared widths added at
collectors cost at most \(O(D^{-4\ell}T^2)\) per cell, even if all
end at the same collector. Additive rounding costs from constantly
many chunks, with \(D\) fixed, obey the same bounds for sufficiently
large \(T\). Combining the three charges using
\((x+y+z)^2\le3(x^2+y^2+z^2)\) proves the squared-arity bound, including
the original external widths. The collector count and Cauchy--Schwarz give
\[
 \sum_{\rm frozen}t_i\le
 \left(O\!\left(\frac W{DT}\right)\,O(WT)\right)^{1/2}
 =O(W/\sqrt D).
\]

\paragraph{Core stations.}
To bound the total arity at the chunk tables, first sum their port widths
over all chunks. Untouched auxiliary insertions have total coded width
\(O(DT(1+\log D))\) per cell by Section~\ref{subsec:plant-private}.
The seams have at most \(g\) ports per cell of total width \(O_g(DT)\)
after summing over chunks in Equation~\eqref{eq:local-chunk-size}.
A cell also has \(O_g(D^2)\) residual private legs. Each costs
\(O_g(DT/\sqrt F)+O(M)\) across all chunks by
Equations~\eqref{eq:plant-widths} and~\eqref{eq:local-chunk-size},
so all residual legs together cost \(O(T)\) per cell for large \(T\).
Legs paired with collectors contribute at most the global frozen arity.
The ordinary wires from motion, the selector wires, and any scalar
tables contribute \(O(T)\) per cell. Thus core arity before controlling the host
degree is \(O(W(1+\log D))\). Untouched auxiliaries have individual
arity \(O(T/a)\), the residues have individual arity
\(O(DT/\sqrt F)\), and the chunk tables and routing wires have each been
made smaller than \(D^{-\ell}T\) by a fixed margin. These and the
selector checks exhaust the core tables.

To build a common host, reserve slots in each cell for each kind of
table and bundle that can occur simultaneously in a single summand.
This requires \(D^{O(1)}\) slots per cell. The block size changes table
entries and selector lengths, while the same slots serve every block
type. Combine scalar factors in their cell. Initially allow all
connections between slots in the same cell or in neighboring cells.
Split any direct connection between collectors with a core identity table
on the thin wire carried along their paths. For each core station, attach a fixed balanced binary tree with one
leaf per possible bundle. Its root receives the table's ordered ports,
and each internal vertex separates its list of bits into the sublists
destined for its two children. The corresponding equality tables pass
each bit along a unique path of length \(O(1+\log D)\) and no internal arity exceeds
twice the original station arity. This proves the asserted degree and
total core bounds. Finally take \(c\) larger than the exponents in these
counts by a margin depending on \(\ell\). Padding a nonempty bundle to
\(\lceil T/D^c\rceil\), or making this minimum charge for any one
station, then uses at most \(O(D^{-2\ell}T)\) extra bits per cell.
Unique zero padding preserves the tensor. It also absorbs any bounded
number per cell of weight-only tables in these counts. This proves the first three
conclusions of Lemma~\ref{lem:local}.

\paragraph{Indices of the outer sum.}
Encode each bundle by its endpoint slots, its length, and its position
in the ordered ports at its ends. The input label already includes any local permutation of bits. No individual exterior address is
needed when a bundle is carried as an ordered list. The branching data
at one occupied cell comprise the initial chunk histograms, the motion
and fusion choices (positions, energies, interpolation points and
orders), and the final chunk splittings. Their number of entries is
bounded for fixed \(D\), including the size of the alphabet of the
histograms. Each integer ranges over \(T^{O_D(1)}\). A field point
uses \(O_D(\log T)\) bits. All private selector bits are wires inside
the summand and are excluded from this count. Zero unused entries
give a unique padding of a term index.

Use the supplied coloring of the \((4D^5+20)\)-th power of the host.
Greedy coloring gives such a coloring with \(O_D(1)\) colors for
fixed \(g\). Start motions successively by color and station slot. Concurrent motions affect disjoint neighborhoods,
since each has at most \(D^5\) crossings. A complete padded schedule
therefore has length bounded in terms of \(D,g\) only. In each step
the cells may send their local descriptions to their neighbors. A
merger concatenates the descriptions of its two input tables and the
operation parameters. A crossing adds the named interpolation, the
fusion, and the new endpoints of its thin wire. The history of any
resulting table draws on a bounded number (depending on \(D\)) of
these records. Thus a label of length \(O_D(A+\log(T+1))\) specifies
each entry as a function of its ports through the fixed dictionary.
The dictionary may use arbitrary interpolation weights, restriction maps
and internal contractions. Generating its labels requires no
evaluation of its entries.

It remains to check that the short computations never require inspecting
the selector bits. Counts in a chunk histogram fix the energies and
coded widths. Range and total checks admit every realizable histogram;
the selector check gives zero whenever an admitted histogram is
unrealizable. Interpolation may use a fixed ordered collection of all
regular field points and weights indexed by the local pole bounds.
The scheduler only checks avoidance of marks, the ranges of the indices,
and the rules of fusion and motion. Zero coefficients may be retained.
The checks and length estimates use arithmetic on
\(O_D(\log(T+1))\)-bit integers or field coordinates, including division
with remainder and rounded square roots. The scheduler copies old labels
as strings in fixed padded slots. At each cell the number of records and
operations on them is bounded independently of \(T\), giving
\(O_D(A+(1+\log T)^{O(1)})\) bit operations and bits communicated.
This proves the final conclusion of Lemma~\ref{lem:local}.
